\subsection{Two-party Product Sampling}
\label{sec:prod}
We next consider the problem of two-party sampling from a product
distribution.  Specifically, given $n$-dimensional vectors
$\vec{w_1}=(w_{1,1},\ldots,w_{1,n})$ and
$\vec{w_2}=(w_{2,1},\ldots,w_{2,n})$ as the private inputs from $P_1$
and $P_2$ respectively, we wish to sample from the distribution $\Dp$
defined by

\[
\Pr[\Dp(\vec{w}_1,\vec{w}_2)=i] = 
\frac{w_{1,i}\cdot w_{2,i}}{\sum_{j=1}^n w_{1,j} \cdot w_{2,j}} = 
\frac{w_{1,i} \cdot w_{2,i}}{\ipw}
\]

Of course, if $\ipw = 0$, the probability space is not well-defined, and
in this case, we require the protocol to simply output $\bot$. 

%\paragraph{Assumptions.}
As before, we assume that all weights in $\vec{w}_1$ and $\vec{w}_2$ are
non-negative.

\iffalse
We also assume, without loss of generality, that $$\|
\vec{w}_1 \|_1 = \| \vec{w}_2 \|_1 = 1;$$ otherwise, the parties can
start the protocol by dividing $\vec{w}_b$ by $\|\vec{w}_b\|_1$, since
we have 

\[
\Pr\left[\Dp\left(\frac{\vec{w}_1}{\|\vec{w}_1\|_1},\frac{\vec{w}_2}{\|\vec{w}_2\|_1}\right)=i\right] = 
\frac{ \frac{w_{1,i}}{\|\vec{w}_1\|_1} \cdot \frac{w_{2,i}}{\|\vec{w}_2\|_1}}{
  \ip{\frac{\vec{w}_1}{\|\vec{w}_1\|_1}}{\frac{\vec{w}_2}{\|\vec{w}_2\|_1}
  }} = \Pr[\Dp(\vec{w}_1,\vec{w}_2)=i].
\]

\paragraph{Communication efficiency and modeling the leakage.}
Our goal is to find a protocol for two-party sampling with sublinear (in
$n$) communication.  However, unlike the case for $L_1$ and $L_2$ sampling, we
show that this goal is actually impossible (see
Section~\ref{sec:prodLB}). Roughly speaking, if parties are allowed to
have arbitrary input vectors, then a sublinear communication solution to
product sampling implies a sublinear communication solution to the
disjointness problem, which is known to be impossible.

Thus, to achieve our original goal we consider restrictions on the
parties' inputs.  In particular, we consider a setting where the inputs
of the two parties are sufficiently correlated, i.e. $\ipw =
\omega\left(\frac{1}{n}\right)$.  Under this assumption, we show a
protocol achieving sublinear communication for product sampling (see
Section~\ref{sec:prodip}).  

As the reader may observe, the performance of our protocol depends on
properties of the parties' inputs.  Thus, just by observing the runtime
of the protocol, an adversary may infer information about the other
party's input.  We explicitly model and bound this leakage to show that
at most the inner product of the two input vectors is leaked.
\fi


\paragraph{Ideal functionality.}
We now define an ideal functionality $\Fp$ for two-party product
sampling. This functionality is parametrized by a function $\flk$
capturing the leakage that the functionality gives to the adversary.  

\bbox
{
\begin{center}
 $\Fp$:  Ideal functionality for two-party product sampling
\end{center}

The functionality has the following parameters:
\BI
\item $n \in \NN$. The dimension of the input weight vectors $\vec{w}_1$ and $\vec{w}_2$.
\item A function $\flk$ describing the leakage.

%\item A leakage function $\ell$ specifies what kind of information the
%functionality leaks about the private inputs of parties $P_1$ and $P_2$. 
\EI

The functionality proceeds as follows:
\BEN
\item Receive inputs $\vec w_1$ and $\vec w_2$ from $P_1$ and $P_2$
respectively. 
\item Compute $\mathsf{leak}=\flk(\vec{w}_1,\vec{w}_2)$
\item If $\ipw=0$, send $\mathsf{leak}$ to the adversary and $\bot$ to
$P_1$ and $P_2$. 
\item Otherwise, sample $i$ with probability $\frac{w_{1,i} \cdot
w_{2,i}}{\ip{\vec{w}_1}{\vec{w}_2}}$, send $\mathsf{leak}$ to the
adversary, and send $i$ to $P_1$ and $P_2$.
\EEN 
}

\medskip


\subsubsection{Impossibility of sublinear product sampling}\label{sec:prodLB}
Our goal is to find a protocol for two-party sampling with sublinear (in
$n$) communication.  However, unlike the case for $L_1$ sampling, we
show that this goal is actually impossible. Roughly speaking, if parties
are allowed to have arbitrary input vectors, then a sublinear
communication solution to product sampling implies a sublinear
communication solution to the disjointness problem, which is known to be
impossible.

For our impossibility result, we first define the two-party disjointness problem.

\paragraph{Disjointness problem.}
The disjointness problem checks if two input sets $S$ and $T$ are
disjoint (i.e., $S \cap T = \emptyset$). Specifically, we consider 
a function $\mathsf{DISJ}^n: \zo^n \times \zo^n \to \zo$ defined as:
$$ \mathsf{DISJ}^n(v_S, v_T) = \left \{ 
    \begin{array}{l} 1 \mbox{ if } \bra v_S,  v_T \ket = 0 \\
    0 \mbox{ otherwise } 
  \end{array} \right .$$ 
In the above, $v_S$ and $v_T$ are the characteristic vectors of $S$ and
$T$ respectively. The communication complexity of the solution to the
disjointness problem is known to have a linear lowerbound, as shown in
the following Theorem:

\begin{theorem}[\cite{tcs:Razborov92,siamdm:KalSch92,jcss:BJKS04}]
For any (even non-private) two-party protocol $\Pi$ where each party
holds $v_S$ and $v_T$ respectively, if $\Pi$ computes $\mathsf{DISJ}^n(v_S,
v_T)$ correctly with probability at least 2/3, the communication
complexity of $\Pi$ is $\Theta(n)$.
\end{theorem}

\paragraph{Our impossibility result.}
We first observe that a simple reduction from Disjointness gives us
that is impossible to achieve sublinear product sampling. Specifically,
disjointness can be directly learned from whether the product sampling
protocol outputs $\bot$ or not. 

Our impossibility result is stronger. We show that it is impossible to
achieve sublinear product sampling \emph{even when the product sampling
protocol is executed with input vectors $\w_1$ and $\w_2$ in which all coordinates are bounded away from $0$}, which in particular guarantees that $\langle \vec{w}_1, \vec{w}_2 \rangle$ is bounded away from $0$. 

%
Before stating a formal theorem below, for $0 < \gamma < 1$, we first
define $\gamma$-heaviness;~we say that a vector $\vec{w}$ is
\emph{$\gamma$-heavy} when each coordinate of $\vec{w}$ is a number
contained in $[\gamma,1]$. 

\begin{theorem} \label{th:lower_bound}
Let $\vec{w}_1$ and $\vec{w}_2$ be 
 $\gamma$-heavy vectors of length $n$, each respectively held by $P_1$
 and $P_2$.   Assume there exists a two-party protocol $\Pp$ for the
 product sampling from $\w_1$ and $\w_2$, with communication at most $C
 := C(n, \gamma)$.  
 
 Then, for any $\gamma \le 1/2n$, there exists a constant $\rho$ and a
 probabilistic protocol computing $\mathsf{DISJ}^n$ correctly with
 probability at least 2/3 that has communication at most $\log(n) + 1 +
 \rho \cdot (C + 1)$.  
 \end{theorem}


\subsubsubsection{Proof of Theorem~\ref{th:lower_bound}} 
We construct a protocol computing $\mathsf{DISJ}^n$ by taking advantage
of $\Pp$ as follows: 

\paragraph{The protocol for $\mathsf{DISJ}^n$}

\noindent
Parties $\mathbf{A}$ and $\mathbf{B}$ each get as input
a vector $\tilde{\vec{a}}$, $\tilde{\vec{b}} \in \{0,1\}^n$.
The goal is to output $1$ if the vectors are ``disjoint''
and $0$ otherwise.

\begin{description}
\item{Edge Case:}
If one of the parties' inputs has Hamming weight $0$, then they output
$1$ and send $1$ to the other party.  From now on, we assume that the
Hamming weight of each party's input is at least $1$.

\item{Preamble:}
We call the party with the lower Hamming weight input the
\emph{designated party}. To determine this, $\mathbf{A}$ sends to
$\mathbf{B}$ the Hamming weight of its input vector $\tilde{\vec{a}}$.
If $\mathbf{B}$'s input has higher Hamming weight, it sends back the bit
$1$ to $\mathbf{A}$; otherwise it sends $0$. 

\item{Input Transformation:}
Let $g_\gamma: \zo \to \RR$ be a boosting function defined as
$g_\gamma(0) = \gamma$ and $g_\gamma(1) = 1$. 
%
Each party $\mathbf{A}$, $\mathbf{B}$ locally transforms their input
vector $\tilde{\vec{a}}$, $\tilde{\vec{b}}$ to $\vec{a}$, $\vec{b}$ by
applying the boosting function in order to ensure $\gamma$-heaviness.
That is, for $i \in [n]$, set $a_i = g_\gamma(\tilde{a}_i)$ and $b_i =
g_\gamma(\tilde{b}_i)$. 

\item{Sampling Protocol:}
The parties run the sampling protocol $\Pp(\vec{a},\vec{b})$ and both receive some output
$i^*$.
\item{Output Computation:}
The \emph{designated party} checks the $i^*$th bit of its input by which
we denote $x$ (i.e., $x = \tilde a_{i^*}$ or $x = \tilde b_{i^*}$
depending on which party is the designated party). It sends $1-x$ to the
other party. Both parties output $1-x$. 
\end{description}

\noindent
The following lemmas give the completeness and soundness of the protocol.

\begin{lemma} \label{clm:completeness} If $\tilde{\vec{a}}$,
$\tilde{\vec{b}}$ are disjoint, then the parties both output $1$ with
probability at least $\frac{1}{2 + n \cdot \gamma}$.  
\end{lemma}

\begin{lemma}  \label{clm:soundness}
If $\tilde{\vec{a}}$, $\tilde{\vec{b}}$ are not disjoint, then the
parties both output $1$ with probability at most $1-\frac{1}{1 + n \cdot
\gamma}$.  
\end{lemma}

Before we prove the lemmas, we briefly describe how we can use these
lemmas to achieve a protocol that correctly computes $\mathsf{DISJ}$
with probability at least $2/3$. 
%
Note that we can get a gap by setting $\gamma = \frac{1}{2n}$.  In other
words, parties output $1$ when disjoint with probability at least
$\frac{2}{5}$.  Parties output $1$ when not disjoint with probability at
most $\frac{1}{3}$. Since we have a constant gap between completeness
and soundness, this can be amplified to $2/3$ and $1/3$ by running the
protocol a constant number of times.

\paragraph{Remarks.}
We would like to characterize the sublinearity condition for product
sampling protocols using the normalized input vectors.  We can do this
since without loss of generality we can assume that input vectors to the
product sampling protocols are normalized; in particular, for any
(non-normalized) vectors $\w_1$ and $\w_2$, we have 

\[
\Pr\left[\Dp\left(\frac{\vec{w}_1}{\|\vec{w}_1\|_1},\frac{\vec{w}_2}{\|\vec{w}_2\|_1}\right)=i\right]
= 
\frac{ \frac{w_{1,i}}{\|\vec{w}_1\|_1} \cdot
\frac{w_{2,i}}{\|\vec{w}_2\|_1}}{
  \ip{\frac{\vec{w}_1}{\|\vec{w}_1\|_1}}{\frac{\vec{w}_2}{\|\vec{w}_2\|_1}
    }} = \Pr[\Dp(\vec{w}_1,\vec{w}_2)=i].
\]

Specifically, we show below that the impossibility theorem implies that
in order to achieve sublinear communication complexity for product
sampling, we would need, at the minimum, a promise on the inputs that
guarantees that $$\langle \vec{w}_1, \vec{w}_2 \rangle \in
\Omega(1/n^2),$$ when $\w_1, \w_2$ are normalized vectors. 

To do this, first note that the theorem implies that sublinear
communication product sampling needs to have $\gamma \in \Omega(1/n)$. 
%
Now, in the proof, any non-disjoint binary vectors $\tilde{\vec{a}}$,
$\tilde{\vec{b}}$ to the $\mathsf{DISJ}$ problem has $\langle
\tilde{\vec{a}}$, $\tilde{\vec{b}} \rangle \ge 1$, and these vectors are
transformed to $g_\gamma(\tilde{\vec{a}})$ and
$g_\gamma(\tilde{\vec{b}})$.
%
Let $\w_1$ and $\w_2$ be the normalized vectors
$g_\gamma(\tilde{\vec{a}})$ and $g_\gamma(\tilde{\vec{b}})$; that is,
$\w_1 := g_\gamma(\tilde{\vec{a}}) / \| g_\gamma(\tilde{\vec{a}}) \|_1$
and $\w_2 =  g_\gamma(\tilde{\vec{b}}) / \| g_\gamma(\tilde{\vec{b}})
\|_1$. Since each entry of $g_\gamma(\tilde{\vec{a}})$ and
$g_\gamma(\tilde{\vec{a}})$ is at most 1, we have $\|g_\gamma(\tilde{\vec{a}})\|_1 \le n$ and $\|g_\gamma(\tilde{\vec{b}})\|_1 \le n$.
Therefore, we have 
%
$$\ipw \ge \frac {\langle g_\gamma(\vec{a}), g_\gamma(\vec{b}) \rangle } { n \cdot n } \ge \frac 1 {n^2}.$$ 

\begin{proof}[Proof of Lemma~\ref{clm:completeness}]
Assume that $\tilde{\vec{a}}$, $\tilde{\vec{b}}$ are disjoint, and
moreover, assume WLOG that $\mathbf{A}$ is the designated party, and its
input vector has Hamming weight $w$.  Recall that $a_i = g_\gamma(\tilde
a_i)$ and $b_i = g_\gamma(\tilde b_i)$. 
Let 
\begin{eqnarray*} 
&& W_{0,0} :=  \sum_{i : \tilde{a}_i=0, \tilde{b}_i=0} a_i \cdot b_i,
~~~~~~ 
W_{1,0}  :=  \sum_{i : \tilde{a}_i=1, \tilde{b}_i=0} a_i \cdot b_i \\
&& W_{0,1}  :=  \sum_{i : \tilde{a}_i=0, \tilde{b}_i=1} a_i \cdot b_i,
~~~~~~
W_{1,1}  :=  \sum_{i : \tilde{a}_i=1, \tilde{b}_i=1} a_i \cdot b_i \\
\end{eqnarray*}

Note that $W_{0,0} \leq n \cdot \gamma^2$.  Further, $W_{1,1} = 0$,
since the vectors are disjoint, and $W_{1,0} = w \cdot \gamma$ since the
Hamming weight of $\tilde{\vec{a}}$ is exactly $w$.  Additionally, note
that $W_{0,1} \geq W_{1,0}$, since $\mathbf{A}$ is the designated party,
so the Hamming weight of $\tilde{\vec{a}}$ is less than or equal to the
Hamming weight of $\tilde{\vec{b}}$.

Note that when the designated party is $\mathbf{A}$, then the output of
the protocol is $1 - a_{i^*}$. Using the above facts, the probability
of outputting $1$ is 
%
\begin{align*}
\frac{W_{0,0} + W_{0,1}}{W_{1,1} + W_{0,0} + W_{0,1} + W_{1,0}}
&\ge 
\frac{W_{0,1}}{W_{0,0} + W_{0,1} + W_{1,0}}\\
&\geq \frac{W_{0,1}}{n\gamma^2 + 2W_{0,1}}\\
&=\frac{w \cdot \gamma}{n\gamma^2 + 2w \cdot \gamma}\\
&=\frac{w}{n\gamma + 2w}\\
&\geq \frac{1}{n\gamma + 2},
\end{align*}
where the last inequality follows since $w \geq 1$, due to the Edge Case
step of the protocol. 
%\qed
\end{proof}

\begin{proof}[Proof of Lemma~\ref{clm:soundness}]
Assume that $\tilde{\vec{a}}$,
$\tilde{\vec{b}}$ are not disjoint.
As before, consider $W_{0,0}$, $W_{1,0}$, $W_{0,1}$, and $W_{1,1}$. 
Note that $W_{1,1} \geq 1$ since the inputs are not disjoint. We also
have $W_{0,0} + W_{0,1} + W_{1,0} \leq n \cdot \gamma$, since $a_i$ or
$b_i$ is $\gamma$ in these cases. 


Using the above facts, the probability of outputting $0$
is 
\begin{align*}
\frac{W_{1,0} + W_{1,1}}{W_{1,1} + W_{0,0} + W_{0,1} + W_{1,0}}
&\geq
\frac{W_{1,1}}{W_{1,1} + n \cdot \gamma}\\
&\ge \frac{1}{1 + n \cdot \gamma}.
\end{align*} 
% \qed
\end{proof}


\subsubsection{Product sampling while leaking at most the inner product}
\label{sec:prodip}

\paragraph{Assumptions.}
As before, we assume that all weights in $\vec{w}_1$ and $\vec{w}_2$ are
non-negative. As discussed in the previous subsection, we also assume,
without loss of generality, that $$\| \vec{w}_1 \|_1 = \| \vec{w}_2 \|_1
= 1.$$ 

\paragraph{Overview.}
We now show that the impossibility result of Section~\ref{sec:prodLB}
can be bypassed if we make some assumptions on the inputs.
Specifically, if we restrict ourselves to the case when
$\ipw=\omega\left(\frac{\log n}{n}\right)$, then we can achieve a sublinear
communication protocol for product sampling on inputs $\vec{w}_1,
\vec{w}_2$\footnote{Regarding $\ipw$, there is a gap between the
lowerbound result (i.e., $\Omega(\frac{1}{n^2})$) and our construction
(i.e., $\omega(\frac{\log n}{n})$). Resolving the gap is left as an
interesting open problem.}.
Of course, by observing that the protocol uses sub-linear
communication, due to our lower-bound, both parties will learn that such
a promise on the inputs is satisfied;~the lower bound implies that some
leakage about the inputs is necessary.  
%
In our protocol, we show that the information leaked is at most the
inner product $\ipw$.
(Formally, we set $\flk(\w_1, \w_2) =\ipw$.)  Interestingly, we show that
this is the case even though our protocol does not, and
cannot,\footnote{This can be shown by a simple modification of the lower
bound proof from Section~\ref{sec:prodLB}.} actually compute $\ipw$.


\paragraph{Product sampling protocol.}
Roughly, the protocol works as follows.  The protocol proceeds in rounds
where in each round $P_1$ and $P_2$ use the oblivious $L_1$ sampling
with a single input vector ($\F_\oslone$) to produce two secret-shared
sampled indices, one from $P_1$'s input vector, and one from $P_2$'s
input vector.  The parties then run a secure 2-PC protocol to securely
compare these values, and if they are equal, output the sampled index.
If the two sampled indices are not equal, the parties move to the next
round.

We describe a private two-party protocol for product sampling leaking at
most the inner product (see Protocol~\ref{prot:prodsampapx}). This
protocol is in the $\{\F_\oslone, \F_{2PC} \}$-hybrid model.

\floatname{algorithm}{Protocol}
\begin{algorithm}[htbp]
\textbf{Inputs:} Party $P_b$ has input $\vec{w}_b$ of length $n$.   

\begin{enumerate}
    \item Invoke the $\F_\oslone$ ideal functionality with $P_1$ as the
    sender with input $\vec{w}_1$ and $P_2$ as the receiver.  Let
    $i_{1,1}$ and $i_{1,2}$ be the output from the ideal functionality
    to $P_1$ and $P_2$ respectively. 
    
    \item Invoke the $\F_\oslone$ ideal functionality with $P_2$ as the
    sender with input $\vec{w}_2$ and $P_1$ as the receiver.  Let
    $i_{2,1}$ and $i_{2,2}$ be the output from the ideal functionality
    to $P_1$ and $P_2$ respectively.
    
    \item Invoke the $\F_{2PC}$ ideal functionality with the following circuit: 
      \BEN
      \item[] \hspace{-2em} Input: $(i_{1,j}, i_{2,j})$ for $j = 1, 2$. 
      \item Let $i_1 = i_{1,1} \xor i_{1,2}$, $i_2 = i_{2,1} \xor i_{2,2}$. 
      \item If $i_1$ is equal to $i_2$, output $i_1$ to both $P_1$ and $P_2$. Otherwise,
      output $\bot$.
      \EEN
  
    \item If the output from the ideal functionality is $\bot$, go back to Step 1.
    Otherwise, output whatever $\F_{2PC}$ outputs. 
\end{enumerate}
\textbf{Output:} Both parties output the sampled value $i$. 

\caption{Product sampling ($\Pp^{IP}$) in the $\{\F_\oslone, \F_{2PC} \}$-hybrid.
\label{prot:prodsampapx}
}
\end{algorithm}


% A value is sampled by each party (without the party knowing the value sampled).
% An equality check is performed.
% If it fails, the parties learn that equality fails.
% If it succeeds, the parties learn the output.

% The parties re-run the above procedure until it succeeds. They output whatever was outputted in the succeeding run.

\paragraph{Security.}
We will prove the following theorem.
\begin{theorem}\label{th:prodIP}
Protocol $\Pp^{IP}$ securely realizes $\Fp$ with leakage $\flk(\w_1,
\w_2) = \ip{\vec{w}_1}{\vec{w}_2}$ in the $\{\F_\oslone, \F_{2PC}
\}$-hybrid model with semi-honest security.
\end{theorem}

\begin{proof}
We describe the simulator $\Sim$ in the $\{\F_\oslone, \F_{2PC} \}$-hybrid model
for the case that Party 1 is corrupted.
The simulator and proof of security are analogous in the case that Party 2 is corrupted.

$\Sim$ receives as input
$\vec{w}_1$, the output $i^*$, and $\ip{\vec{w}_1}{\vec{w}_2}$.
$\Sim$ samples $r^*$ from a geometric distribution with success probability $p = \langle
\vec{w}_1, \vec{w}_2 \rangle$.

$\Sim$ invokes Party 1 on input $\vec{w}_1$.
For $i \in [r^*-1]$, 
Party 1 sends its input to the first
invocation of 
$\F_\oslone$ and $\Sim$
returns to it a random value in $\mathbb{Z}_n$.
Party 1 sends its input to the second
invocation of 
$\F_\oslone$ and $\Sim$
returns to it a random value in $\mathbb{Z}_n$.
Party 1 sends its input to the 
$\F_{2PC}$ functionality and $\Sim$
returns to it $\bot$.
For $i = r^*$,
Party 1 sends its input to the first
invocation of 
$\F_\oslone$ and $\Sim$
returns to it a random value in $\mathbb{Z}_n$.
Party 1 sends its input to the second
invocation of 
$\F_\oslone$ and $\Sim$
returns to it a random value in $\mathbb{Z}_n$.
Party 1 sends its input to the 
$\F_{2PC}$ functionality and $\Sim$
returns to it $i^*$.

It is clear that the view of Party 1
is identical in the ideal and real world,
assuming that $\Sim$
samples the first succeeding round,
$r^*$, from the correct distribution.
In the following, we argue that this is 
indeed the case.

\iffalse
Note that each round of the protocol consists of two runs of the
oblivious sampling protocol. Thus, it can be simulated by having the
simulator playing as $F_\oslone$.

The only thing left to simulate is the number of rounds, $r^*$. The simulator receives
from the ideal functionality $\Fp$ the sampled output $i^*$.  We show
that given $\fl=\ip{\vec{w}_1}{\vec{w}_2}$ and the output $i^*$, the
simulator can sample the first succeeding round $r^*$ from the exact
correct distribution.
\fi

First, note that on any given round, we have 
$$p_c := \Pr[\mbox{collision}] = \sum_i \Pr[i_i = i \wedge i_2 = i] = \sum_i w_{1,i} \cdot w_{2,i} = \ipw.$$ 


Let ${\sf FirstSuccess}(r)$ denote an event in which the protocol
succeeds for the first time on the $r$-th round.  Now, for $r \in
\mathbb{N}$, we have 
%
\begin{align*}
&\Pr[{\sf FirstSuccess}(r) \mbox{ AND the output is } i^*]\\
&=
\Pr[\mbox{no collision in first } r-1 \mbox{ rounds}] \cdot \Pr[i_1 = i^* \wedge i_2 = i^* \mbox{ on the $r$th round}]\\
&= (1-p_c)^{r-1} \cdot \Pr[i_1 = i^* \wedge i_2 = i^*] 
\end{align*}
%
Now, the probability that the protocol eventually outputs $i^*$ is:
\begin{align*}
&\Pr[\mbox{protocol eventually outputs } i^* \mbox{ after some number of rounds}]\\
&= \sum_{j=1}^\infty \Pr[{\sf FirstSuccess}(j) \mbox{ AND the output is } i^*] \\ &= \Pr[i_1 = i^* \wedge i_2 = i^*]
\sum_{j=1}^\infty (1-p_c)^{j-1}
= \Pr[i_1 = i^* \wedge i_2 = i^*] \cdot \frac{1}{p_c}.
\end{align*}
Thus, the probability of ${\sf FirstSuccess}(r)$ conditioned on the
output being $i^*$ is:
\begin{align*}
&\Pr[{\sf FirstSuccess}(r) | \mbox{ the output is $i^*$}]\\
&= \frac{\Pr[{\sf FirstSuccess}(r) \mbox{ AND the output is }
i^*]}{\Pr[\mbox{protocol eventually outputs } i^* \mbox{ after some number of rounds}]}\\
&= \frac{(\Pr[i_1 = i^* \wedge i_2 = i^*]) \cdot (1-p_c)^{r-1}}{\Pr[i_1
= i^* \wedge i_2 = i^*] \cdot \frac{1}{p_c}}\\
&= p_c \cdot (1-p_c)^{r-1}.
\end{align*}

The above is exactly the probability of the number of Bernoulli trials
(with probability $p_c = \langle \vec{w}_1, \vec{w}_2 \rangle$) needed to
get one success.  Sampling the number of rounds is therefore equivalent
to sampling the random variable corresponding to the number of rounds
from a geometric distribution with success probability $p_c = \langle
\vec{w}_1, \vec{w}_2 \rangle$,
which is exactly what $\Sim$ does. 
% \qed
%This can be efficiently simulated given
%just the dot product $p = \langle \vec{w}_1, \vec{w}_2 \rangle$. 
\end{proof}

\paragraph{Performance.}
As shown above, the number of rounds $r$ needed by this protocol is
distributed as the number of Bernoulli trials (with probability $p =
\ipw$) needed to get one success.  Thus, the expected number of rounds
is $r=\frac{1}{\ipw}$.  In each round, the communication consists of a
secure 2-PC of equality on $O(\log n)$-bit inputs, which can be done in
$O(\log n)$ communication and $O(1)$ rounds.  Thus, in total, this
protocol has expected communication $O(\frac{\log n}{\ipw})$ and
$O(\frac{1}{\ipw})$ rounds.  This communication is sublinear in $n$ when
$\ipw = \omega\left(\frac{\log n}{n}\right)$.\medskip

\noindent \textbf{Trading efficiency for privacy.}  In the proof above, the simulator requires the value of $\ipw$, which is not revealed by the output.  However, a slight modification to the protocol allows us to remove this leakage at the cost of additional, though still sub-linear, communication.  
Instead of terminating the protocol the first time there is a collision in the $L_1$ samples, we can pad the communication cost by making $O(\frac{n}{\log n})$ 
calls to $\F_{\oslone}$.  Under the promise of $\ipw = \omega(\frac{\log n}{n})$, this ensures a collision in 
the outputs (with all but negligible probability).  The parties can then use $O(\frac{n}{\log n})$ communication to obliviously find and output the collision, without revealing the index, and avoiding the leakage of $\ipw$.  

Generalizing this idea, we arrive at a set of similar protocol modifications that support a continuous set of tradeoffs:~instead of choosing between  leaking $\ipw$ to the simulator, or padding to the maximum communication, we 
can choose to leak some lower bound on $\ipw$, and modify the protocol to make a proportionate number of calls to $\F_{\oslone}$, search (obliviously) for a collision, and repeat if necessary.  

Without a full proof, we provide some intuition for the fact that this tradeoff between leakage and communication is inherent.  We can do that by generalizing the statement of Theorem \ref{th:lower_bound}.  We first modify the definition of $\gamma$-heavy defined previously:~for any $t(n) = O(n)$, we say that a vector $\w$ of length $n$ is $\gamma_{t,n}$-heavy if each of the $t := t(n)$ coordinates of $\w$ is a number contained in $[\gamma, 1]$. 
In particular, we now allow $t(n) = o(n)$.  Then, with a small modification to the reduction, we can prove that if $\w_1$ and $\w_2$ are $\gamma_{t,n}$-heavy, and if there exists a protocol $\Pp$ for product sampling with communication at most $C := C(n, \gamma)$, then there exists a protocol for computing $\mathsf{DISJ}^t$ with communication $\log(n) + O(C)$.  In the modified reduction, the parties simply increase the weights of the $t$ input slots (as before), and append $n-t$ entries containing 0 at the end.  
Since we know that $\mathsf{DISJ}^t$ requires $O(t)$ communication, the implication is that we have increasingly weaker communication bounds as we are provided increasingly strong promises on the inner product.  
Conversely, for a certain set of input vectors, observing the communication of the sampling protocol gives you a bound on the inner product of the inputs.  The less communication observed, the tighter that bound, and the greater the leakage.  






